\begin{exercise}[subtitle=Atombau, points=0.5, Punkt=0.5] 
Welche Teilchen befinden sich im Atomkern?
\begin{choices}(2) %% Anzahl der Spalten
  \task Elektronen
  \task \answer{Protonen und Neutronen}
  \task Protonen und Elektronen
  \task Neutronen und Elektronen
\end{choices}
\begin{solution}
Die korrekte Antwort ist \textbf{\GetExerciseProperty{answer}}.
\end{solution}
\end{exercise}

\begin{exercise}[subtitle=Periodensystem, points=0.5, Punkt=0.5] %% subtitle= und points= sind optional
Welches Element hat die Ordnungszahl 11?
\begin{choices}(2)
  \task Lithium
  \task Kalium
  \task \answerNatrium
  \task Magnesium
\end{choices}
\begin{solution}
Die korrekte Antwort ist \textbf{\GetExerciseProperty{answer}}.
\end{solution}
\end{exercise}

\begin{exercise}[subtitle=Ionische Bindung, points=0.5, Punkt=0.5]
Welche der folgenden Verbindungen ist hauptsächlich ionisch gebunden?
\begin{choices}(2)
  \task \(\mathrm{CO_2}\)
  \task \(\mathrm{H_2O}\)
  \task \(\mathrm{NH_3}\)
  \task \answer{\(\mathrm{NaCl}\)}
\end{choices}
\begin{solution}
Die korrekte Antwort ist \textbf{\GetExerciseProperty{answer}}.
\end{solution}
\end{exercise}

\begin{exercise}[subtitle=Atommassen, points=0.5, Punkt=0.5]
Die molare Masse von Sauerstoff \(\mathrm{O_2}\) beträgt ungefähr:
\begin{choices}(2)
  \task 8 g/mol
  \task \answer{32 g/mol}
  \task 16 g/mol
  \task 64 g/mol
\end{choices}
\begin{solution}
Die korrekte Antwort ist \textbf{\GetExerciseProperty{answer}}.
\end{solution}
\end{exercise}

\begin{exercise}[subtitle=Avogadro-Zahl, points=0.5, Punkt=0.5]
Die Avogadro-Zahl beträgt:
\begin{choices}(2)
  \task \(\approx 6{,}02 \times 10^{20}\)
  \task \answer{\(\approx 6{,}02 \times 10^{23}\)}
  \task \(\approx 6{,}02 \times 10^{26}\)
  \task \(\approx 6{,}02 \times 10^{30}\)
\end{choices}
\begin{solution}
Die korrekte Antwort ist \textbf{\GetExerciseProperty{answer}}.
\end{solution}
\end{exercise}

\begin{exercise}[subtitle=Oxidationszahlen, points=0.5, Punkt=0.5]
Welche Oxidationszahl besitzt Sauerstoff in \(\mathrm{H_2O}\)?
\begin{choices}(2)
  \task +2
  \task 0
  \task \answer{-2}
  \task -1
\end{choices}
\begin{solution}
Die korrekte Antwort ist \textbf{\GetExerciseProperty{answer}}.
\end{solution}
\end{exercise}

\begin{exercise}[subtitle=Säure-Base, points=0.5, Punkt=0.5]
Welche der folgenden Substanzen ist eine starke Säure?
\begin{choices}(2)
  \task \(\mathrm{HF}\)
  \task \(\mathrm{CH_3COOH}\)
  \task \answer{\(\mathrm{HCl}\)}
  \task \(\mathrm{H_2O}\)
\end{choices}
\begin{solution}
Die korrekte Antwort ist \textbf{\GetExerciseProperty{answer}}.
\end{solution}
\end{exercise}

\begin{exercise}[subtitle=Gasgesetz, points=0.5, Punkt=0.5]
Bei konstantem Druck und Stoffmenge ist das Volumen eines idealen Gases proportional zu...
\begin{choices}(2)
  \task dem Kehrwert der Temperatur.
  \task \answer{der Temperatur.}
  \task dem Druck.
  \task der Molmasse.
\end{choices}
\begin{solution}
Die korrekte Antwort ist \textbf{\GetExerciseProperty{answer}}.
\end{solution}
\end{exercise}

\begin{exercise}[subtitle=Chemische Gleichung, points=0.5, Punkt=0.5]
Welche stöchiometrische Koeffizienten balancieren die Gleichung \(\mathrm{C_3H_8 + O_2 \rightarrow CO_2 + H_2O}\)?
\begin{choices}(2)
  \task \(\mathrm{1:3:3:4}\)
  \task \(\mathrm{2:7:6:8}\)
  \task \answer{\(\mathrm{1:5:3:4}\)}
  \task \(\mathrm{2:5:4:4}\)
\end{choices}
\begin{solution}
Die korrekte Antwort ist \textbf{\GetExerciseProperty{answer}}.
\end{solution}
\end{exercise}

\begin{exercise}[subtitle=Elektronegativität, points=0.5, Punkt=0.5]
Welches Element hat die höchste Elektronegativität nach der Pauling-Skala?
\begin{choices}(2)
  \task Sauerstoff
  \task Stickstoff
  \task Chlor
  \task \answer{Fluor}
\end{choices}
\begin{solution}
Die korrekte Antwort ist \textbf{\GetExerciseProperty{answer}}.
\end{solution}
\end{exercise}

\begin{exercise}[subtitle=Lewis-Struktur, points=0.5, Punkt=0.5]
Wie viele bindende Elektronenpaare hat Methan (\(\mathrm{CH_4}\)) in seiner Lewis-Struktur?
\begin{choices}(2)
  \task 3
  \task \answer{4}
  \task 5
  \task 6
\end{choices}
\begin{solution}
Die korrekte Antwort ist \textbf{\GetExerciseProperty{answer}}.
\end{solution}
\end{exercise}

\begin{exercise}[subtitle=Massenwirkungsgesetz, points=0.5, Punkt=0.5]
Für die Gleichgewichtsreaktion \(\mathrm{aA + bB \leftrightarrow cC + dD}\) lautet die Gleichgewichtskonstante \(K\):
\begin{choices}(2)
  \task \(\dfrac{[A]^a[B]^b}{[C]^c[D]^d}\)
  \task \answer{\(\dfrac{[C]^c[D]^d}{[A]^a[B]^b}\)}
  \task \(\dfrac{aA + bB}{cC + dD}\)
  \task \(\dfrac{cC + dD}{aA + bB}\)
\end{choices}
\begin{solution}
Die korrekte Antwort ist \textbf{\GetExerciseProperty{answer}}.
\end{solution}
\end{exercise}

\begin{exercise}[subtitle=Redoxreaktion, points=0.5, Punkt=0.5]
In der Reaktion \(\mathrm{Zn + Cu^{2+} \rightarrow Zn^{2+} + Cu}\) fungiert Zink als...
\begin{choices}(2)
  \task Oxidationsmittel
  \task \answer{Reduktionsmittel}
  \task Katalysator
  \task Edelmetall
\end{choices}
\begin{solution}
Die korrekte Antwort ist \textbf{\GetExerciseProperty{answer}}.
\end{solution}
\end{exercise}

\begin{exercise}[subtitle=Begriff der Molalität, points=0.5, Punkt=0.5]
Die Molalität definiert sich als...
\begin{choices}(2)
  \task Anzahl Mol des gelösten Stoffes pro Liter Lösung.
  \task \answer{Anzahl Mol des gelösten Stoffes pro Kilogramm Lösungsmittel.}
  \task Anzahl Mol pro Kubikmeter Gas.
  \task Anzahl Mol des Lösungsmittels pro Kilogramm Lösung.
\end{choices}
\begin{solution}
Die korrekte Antwort ist \textbf{\GetExerciseProperty{answer}}.
\end{solution}
\end{exercise}

\begin{exercise}[subtitle=Neutralisationsreaktion, points=0.5, Punkt=0.5]
Welche Produkte entstehen typischerweise bei der Reaktion einer Säure mit einer Base?
\begin{choices}(2)
  \task \answer{Salz und Wasser}
  \task Gas und Wasser
  \task Säure und Base
  \task Salz und Gas
\end{choices}
\begin{solution}
Die korrekte Antwort ist \textbf{\GetExerciseProperty{answer}}.
\end{solution}
\end{exercise}
